%==========================================================================
% BCC502 -- Metodologia Científica para Ciência da Computação
% Template da tarefa: proposta de pesquisa (contextualização, justificativa
% da hipótese e objetivos)
%
% Regras:
%  1. Preencha \titulotrabalho, \nomealuno e \matricula e o conteúdo dos slides.
%  2. A capa é obrigatória. Entregue de 4 a 6 slides de conteúdo, além da capa.
%     O conteúdo e os títulos dos slides são livres. Os slides abaixo são
%     sugestões: você pode apagar, duplicar, fundir ou reordenar.
%  3. Cada slide deve trazer, no subtítulo (segundo argumento do frame), o
%     campo que endereça: Campo 1 (Contextualização e problema), Campo 2
%     (Justificativa da hipótese) ou Campo 3 (Objetivos). Todos os três campos
%     devem aparecer.
%  4. Não altere as cores, as fontes nem o rodapé.
%  5. Apague todo texto entre colchetes vermelhos antes de entregar.
%  6. Entregue um único PDF, nomeado matricula_nome.pdf.
%==========================================================================
\documentclass[aspectratio=169,11pt]{beamer}

\usepackage[utf8]{inputenc}
\usepackage[T1]{fontenc}
\usepackage[brazil]{babel}
\usepackage{microtype}
\usepackage{csquotes}
\usepackage{booktabs}
\usepackage{array}
\usepackage{tabularx}
\usepackage[most]{tcolorbox}

%--------------------------------------------------------------------------
% DADOS DO ALUNO (preencha)
%--------------------------------------------------------------------------
\newcommand{\titulotrabalho}{Título provisório do trabalho}
\newcommand{\nomealuno}{Nome Completo do Aluno}
\newcommand{\matricula}{00.0.0000}

%--------------------------------------------------------------------------
% Cores e tema (não alterar)
%--------------------------------------------------------------------------
\definecolor{mtext}{RGB}{38, 38, 38}
\definecolor{mgray}{RGB}{120, 120, 120}
\definecolor{mlight}{RGB}{215, 215, 215}
\definecolor{maccent}{HTML}{5D5179}
\definecolor{mplace}{RGB}{176, 58, 46}

\usetheme{default}
\usecolortheme{default}
\setbeamercolor{normal text}{fg=mtext, bg=white}
\setbeamercolor{frametitle}{fg=mtext}
\setbeamercolor{itemize item}{fg=maccent}
\setbeamercolor{itemize subitem}{fg=mgray}

\setbeamerfont{frametitle}{series=\bfseries, size=\large}
\setbeamerfont{captitulo}{series=\bfseries, size=\fontsize{22}{27}}

\setbeamersize{text margin left=1.3cm, text margin right=1.3cm}
\setbeamertemplate{navigation symbols}{}
\setbeamertemplate{itemize items}[circle]
\setlength{\parskip}{0pt}

% Título: texto à esquerda e um filete curto de destaque
\setbeamertemplate{frametitle}{%
  \vskip 0.15cm
  {\usebeamerfont{frametitle}\usebeamercolor[fg]{frametitle}\insertframetitle\hfill{\footnotesize\bfseries\color{maccent}\insertframesubtitle}\par}
  \vskip 0.1cm
  \textcolor{maccent}{\rule{1.1cm}{1.5pt}}
  \vskip -0.2cm
}

% Rodapé discreto
\setbeamertemplate{footline}{%
  \leavevmode%
  \hbox{%
    \begin{beamercolorbox}[wd=0.8\paperwidth, ht=2.5ex, dp=1.2ex, leftskip=1.3cm]{normal text}%
      \scriptsize\color{mgray} BCC502 \textperiodcentered{} \nomealuno{} \textperiodcentered{} Matrícula \matricula
    \end{beamercolorbox}%
    \begin{beamercolorbox}[wd=0.2\paperwidth, ht=2.5ex, dp=1.2ex, rightskip=1.3cm, right]{normal text}%
      \scriptsize\color{mgray}\insertframenumber{}
    \end{beamercolorbox}%
  }%
}

% Bloco sem preenchimento, só com traço lateral
\newtcolorbox{bloco}[1]{%
  enhanced, frame hidden, boxrule=0pt, arc=0pt,
  colback=white, colbacktitle=white,
  borderline west={1.5pt}{0pt}{maccent},
  left=7pt, right=2pt, top=3pt, bottom=3pt,
  before skip=5pt, after skip=2pt,
  before upper={\textcolor{maccent}{\small\bfseries #1}\par\vspace{0.2em}},
}

% Rótulo de coluna ou de lista
\newcommand{\rotulo}[1]{\textcolor{maccent}{\small\bfseries #1}}

% Texto a ser substituído pelo aluno
\newcommand{\preencher}[1]{\textcolor{mplace}{\textit{[#1]}}}

% Tabelas: apenas linhas horizontais finas
\setlength{\heavyrulewidth}{0.8pt}
\setlength{\lightrulewidth}{0.4pt}
\newcolumntype{L}{>{\raggedright\arraybackslash}X}

\begin{document}

%==========================================================================
% CAPA (obrigatória; não conta como slide de conteúdo)
%==========================================================================
\begin{frame}[plain,c]
  \relax
  {\usebeamercolor[fg]{frametitle}\usebeamerfont{captitulo}\titulotrabalho\par}
  \vskip 0.35cm
  \textcolor{maccent}{\rule{1.1cm}{1.5pt}}
  \vskip 0.5cm
  {\large \nomealuno\par}
  \vskip 0.15cm
  {\small\color{mgray} BCC502 \textperiodcentered{} Matrícula \matricula\par}
\end{frame}
\setcounter{framenumber}{0}

%==========================================================================
% SLIDES DE CONTEÚDO (sugestões; conteúdo livre)
% O subtítulo (segundo argumento) indica o campo endereçado.
%==========================================================================

\begin{frame}[t]{Contextualização}{Campo 1 \textperiodcentered{} Contextualização e problema}
  \begin{bloco}{A área e o que se faz hoje}
    \preencher{A área e o objeto em uma ou duas frases acessíveis, e as abordagens existentes mais próximas, com fonte (autor, ano).}
  \end{bloco}

  \begin{bloco}{Onde as abordagens atuais falham}
    \preencher{A limitação concreta, de preferência como custo e benefício, e a evidência de que ela existe.}
  \end{bloco}
\end{frame}

\begin{frame}[t]{O problema de pesquisa}{Campo 1 \textperiodcentered{} Contextualização e problema}
  \begin{bloco}{O problema}
    \preencher{O problema em uma ou duas frases, que um leitor de fora da área consiga repetir.}
  \end{bloco}

  \begin{bloco}{Por que pede pesquisa e não apenas implantação}
    \preencher{Que conhecimento novo está em jogo. Se for propor algo novo, o que existe de mais parecido e em que difere.}
  \end{bloco}
\end{frame}

\begin{frame}[t]{Hipótese}{Campo 2 \textperiodcentered{} Justificativa da hipótese}
  \begin{bloco}{Hipótese}
    \preencher{Uma frase: o que acredito que seja verdade e ainda não está demonstrado.}
  \end{bloco}

  \begin{bloco}{Que limitação ela ataca}
    \preencher{Qual limitação apontada no slide anterior e por que esta abordagem deveria superá-la.}
  \end{bloco}
\end{frame}

\begin{frame}[t]{Por que acredito na hipótese}{Campo 2 \textperiodcentered{} Justificativa da hipótese}
  \begin{bloco}{Evidências}
    \begin{itemize}\small
      \item \rotulo{Trabalhos que apontam viabilidade:} \preencher{autor, ano, o que mostraram}
      \item \rotulo{Dados preliminares meus:} \preencher{o que testei e o que observei}
      \item \rotulo{Estudo de caso:} \preencher{se houver}
    \end{itemize}
  \end{bloco}
\end{frame}

\begin{frame}[t]{Objetivos}{Campo 3 \textperiodcentered{} Objetivos}
  \begin{bloco}{Objetivo geral}
    \preencher{Verbo não trivial, o que se obtém, a medida e o valor de referência com o qual se compara.}
  \end{bloco}

  \begin{bloco}{Objetivos específicos}
    \begin{itemize}\small
      \item \preencher{Resultado verificável que detalha o geral, com verbo}
      \item \preencher{Outro resultado verificável, e não um passo do método}
    \end{itemize}
  \end{bloco}
\end{frame}

\end{document}